\documentclass{amsart}
\usepackage{amsmath,amssymb,amsthm,hyperref}
\usepackage{algorithm}
\usepackage{algpseudocode}
\newtheorem{theorem}{Theorem}
\newtheorem{lemma}{Lemma}
\newtheorem{proposition}{Proposition}
\newtheorem{corollary}{Corollary}
\theoremstyle{definition}
\newtheorem{definition}{Definition}
\newtheorem{remark}{Remark}

\begin{document}

\title{Geodetic homeomorphs of the Hoffman--Singleton graph:
a complete characterization and enumeration framework}
\maketitle

\section{Statement and conventions}

Let $H$ be the Hoffman--Singleton graph: the unique $(7,5)$--Moore graph, which
is $7$--regular on $50$ vertices, has $|E(H)|=175$ edges, girth $5$, diameter
$2$, and is geodetic (between any two vertices there is a unique shortest path).
For a Moore graph of diameter $2$ this is automatic: adjacent vertices have no
common neighbour (no triangles), and non-adjacent vertices have exactly one
common neighbour, so every shortest path is unique. The only structural fact
about $H$ that the arguments below use, beyond geodeticity, is that
\emph{$H$ has girth $5$}.

For a positive-integer assignment $x=(x_e)_{e\in E(H)}\in\mathbb Z_{>0}^{175}$
let $G'=G'(x)$ be the subdivision of $H$ in which each edge $e=uv$ is replaced
by an internally disjoint path of length $x_e$ joining $u$ and $v$. The $50$
original vertices, of degree $7$, are the \emph{branch vertices}; the new
degree-$2$ vertices are \emph{internal}. We must decide for which $x$ the graph
$G'(x)$ is geodetic, enumerate all such $x$, and give an algorithm doing so.

Throughout we regard $H$ as a \emph{weighted} graph in which each edge $e$ has
weight $x_e>0$. An \emph{$H$--path} between branch vertices $u,v$ is a simple
path in $H$ from $u$ to $v$; its \emph{weight} is $\sum_{e\in P}x_e$. We write
$D^x(u,v)$ for the minimum weight of an $H$--path from $u$ to $v$ (the weighted
shortest-path distance), and $d_{G'}$ for graph distance in $G'(x)$.

\begin{lemma}[Structure of $H$]\label{lem:struct}
$H$ is $7$--regular on $50$ vertices with $175$ edges, has girth $5$, diameter
$2$, is triangle-free, and has adjacency spectrum
$\{7^{1},2^{28},(-3)^{21}\}$. It contains exactly $1260$ pentagons ($5$--cycles).
\end{lemma}

\begin{proof}
These are the standard properties of the unique $(7,5)$--Moore graph. The
pentagon count follows from the spectrum: since $H$ is triangle-free, the number
of $5$--cycles is
$\tfrac{1}{10}\operatorname{tr}(A^{5})=\tfrac{1}{10}(7^{5}+28\cdot2^{5}
+21\cdot(-3)^{5})=\tfrac{1}{10}\cdot12600=1260$.
\end{proof}

\section{Distances in $G'(x)$ reduce to weighted $H$--paths}

\begin{lemma}\label{lem:distbranch}
For branch vertices $u,v$, $d_{G'(x)}(u,v)=D^x(u,v)$, and the number of distinct
shortest $u$--$v$ paths in $G'(x)$ equals the number of distinct minimum-weight
$H$--paths from $u$ to $v$. Moreover, every subpath of a minimum-weight
$H$--path is itself the unique minimum-weight $H$--path between its endpoints
whenever \textnormal{(U)} below holds.
\end{lemma}

\begin{proof}
Every internal vertex of $G'(x)$ has degree $2$; thus any walk in $G'(x)$, on
entering a subdivided edge $e=ab$ at one endpoint, must traverse it entirely to
the other endpoint before leaving. Hence every $u$--$v$ walk between branch
vertices is determined by the sequence of branch vertices it visits (an
$H$--walk) and has $G'$--length equal to the total weight of the traversed
edges. A minimum-length such walk repeats no branch vertex (deleting a closed
sub-walk strictly shortens it, all weights being positive), so it is an
$H$--path; this gives the distance formula. Distinct internally disjoint
$u$--$v$ paths of $G'(x)$ project to distinct $H$--paths and conversely each
$H$--path lifts to exactly one $u$--$v$ path, so the counts match. Finally, if
$P$ is a minimum-weight $u$--$v$ $H$--path and $a,b$ lie on $P$, the subpath
$P_{a\to b}$ has weight $D^x(a,b)$: any lighter $a$--$b$ walk would, spliced into
$P$, yield an $a$--$b$ replacement producing a $u$--$v$ walk of weight
$<D^x(u,v)$, impossible.
\end{proof}

For a branch vertex $\alpha$ and an internal vertex $q$ lying on edge $e=ab$ at
$G'$--distance $s$ from $a$ ($0<s<x_e$), the same degree-$2$ rigidity gives
\begin{equation}\label{eq:internaldist}
d_{G'}(\alpha,q)=\min\!\big(d_{G'}(\alpha,a)+s,\ d_{G'}(\alpha,b)+(x_e-s)\big)
=\min\!\big(D^x(\alpha,a)+s,\ D^x(\alpha,b)+(x_e-s)\big),
\end{equation}
because every $\alpha$--$q$ path enters $e$ at $a$ or at $b$ and then proceeds
monotonically along $e$ to $q$. We extend the notation $d_{G'}(\alpha,q)$ to all
vertices $q$ by~\eqref{eq:internaldist} (with $d_{G'}(\alpha,v)=D^x(\alpha,v)$
for branch $v$).

\section{The characterization theorem}

\begin{definition}\label{def:conditions}
For $x\in\mathbb Z_{>0}^{175}$ define:
\begin{itemize}
\item[\textnormal{(U)}] (\emph{branch uniqueness}) for every ordered pair of
distinct branch vertices $u,v$ the minimum-weight $H$--path from $u$ to $v$ is
unique;
\item[\textnormal{(G)}] (\emph{no opposite-direction tie}) for every edge
$e=ab$ with $x_e\ge2$, every internal offset $r\in\{1,\dots,x_e-1\}$, and every
vertex $q$ of $G'(x)$ that does not lie on $e$,
\[
 r+d_{G'}(a,q)\;\neq\;(x_e-r)+d_{G'}(b,q).
\]
\end{itemize}
We single out the special case of \textnormal{(G)} in which $q$ is a
\emph{branch} vertex; call it \textnormal{(I)}.
\end{definition}

The correct condition is \textnormal{(G)}, in which $q$ ranges over \emph{all}
vertices --- branch and \emph{internal}. The weaker branch-only condition
\textnormal{(I)} is necessary but \emph{not} sufficient
(Remark~\ref{rem:counterex}).

\begin{theorem}[Characterization]\label{thm:main}
For $x\in\mathbb Z_{>0}^{175}$, the homeomorph $G'(x)$ is geodetic if and only if
both \textnormal{(U)} and \textnormal{(G)} hold. Both directions are proved in
full.
\end{theorem}

The proof of sufficiency uses one auxiliary lemma, which is the crux of the whole
argument; it closes in closed form the on-edge degenerate case that earlier
treatments left to computation.

\begin{lemma}[No redundant balanced edge]\label{lem:redundant}
Assume \textnormal{(U)}. Call an edge $e=ab$ of $H$ \emph{redundant} if the
minimum-weight $a$--$b$ $H$--path avoids $e$, i.e.\ $R:=D^x(a,b)<x_e$; by
\textnormal{(U)} this minimum path $\pi$ is then unique and edge-disjoint from
$e$. If some redundant edge $e$ has $x_e+R$ \emph{even}, then \textnormal{(G)}
fails.
\end{lemma}

\begin{proof}
Let $e=ab$ be redundant with $x_e+R$ even, and write $m:=\tfrac12(x_e+R)\in
\mathbb Z$. Because $\pi$ avoids the edge $e$, the subgraph $e\cup\pi$ is a cycle
of $H$; its length (number of $H$--edges) is at least
$\operatorname{girth}(H)=5$, so $\pi$ has at least $4$ edges and therefore passes
through an \emph{internal branch vertex} $v\notin\{a,b\}$. Put
\[
 \alpha:=D^x(a,v),\qquad \beta:=D^x(b,v).
\]
Since $\pi$ is the minimum-weight $a$--$b$ path, Lemma~\ref{lem:distbranch}
(subpath part) gives that the two arcs of $\pi$ realize $\alpha$ and $\beta$ and
$\alpha+\beta=R$, with $0<\alpha,\beta<R$ because $v$ is interior to $\pi$.

Set $s:=x_e+\beta-m$. From $0<\beta<R$ and $m=\tfrac12(x_e+R)$ we get
$x_e-m<s<m$; moreover $x_e-m=\tfrac12(x_e-R)>0$ (as $e$ is redundant) and
$m=\tfrac12(x_e+R)<x_e$ (again $R<x_e$). Hence $s$ is an integer with
$1\le s\le x_e-1$: a genuine internal offset of $e$. Now compute, using
$d_{G'}(a,v)=\alpha$ and $d_{G'}(b,v)=\beta$ (Lemma~\ref{lem:distbranch}),
\[
 s+d_{G'}(a,v)=(x_e+\beta-m)+\alpha=x_e+R-m=2m-m=m,
\]
\[
 (x_e-s)+d_{G'}(b,v)=\big(x_e-(x_e+\beta-m)\big)+\beta=m.
\]
Thus for the edge $e$, the offset $r:=s\in\{1,\dots,x_e-1\}$ and the target
$q:=v$ --- a \emph{branch} vertex with $v\neq a,b$, so $v$ does not lie on $e$
--- we have $r+d_{G'}(a,q)=(x_e-r)+d_{G'}(b,q)=m$. This is precisely a violation
of \textnormal{(G)}.
\end{proof}

\begin{proof}[Proof of Theorem~\ref{thm:main}]
\emph{Necessity.} If (U) failed, some branch pair would have two minimum-weight
$H$--paths, which by Lemma~\ref{lem:distbranch} lift to two shortest paths in
$G'(x)$. If (G) failed, fix the witnessing $e=ab$, $r$, $q$ and let $p$ be the
internal vertex of $e$ at distance $r$ from $a$. As $p$ is internal (degree $2$),
every $p$--$q$ path leaves $p$ towards $a$ or towards $b$, so
$d_{G'}(p,q)=\min(r+d_{G'}(a,q),\,(x_e-r)+d_{G'}(b,q))$. The failure of (G) makes
the two arguments equal, so both attain $d_{G'}(p,q)$: there is a shortest
$p$--$q$ path leaving $p$ towards $a$ and another leaving towards $b$. These use
different first edges at $p$, hence are distinct. In either case $G'(x)$ is not
geodetic.

\smallskip
\emph{Sufficiency.} Assume (U) and (G). We first record:

\smallskip
\noindent\textbf{Claim (branch source).} For every branch vertex $\alpha$ and
every vertex $q$, the shortest $\alpha$--$q$ path in $G'(x)$ is unique.

\noindent\emph{Proof of Claim.} If $q$ is branch this is (U) via
Lemma~\ref{lem:distbranch}. If $q$ is internal on $f=cd$ at offset $s$, then by
\eqref{eq:internaldist} $d_{G'}(\alpha,q)=\min(D^x(\alpha,c)+s,\,
D^x(\alpha,d)+(x_f-s))$. By (U) each of $D^x(\alpha,c),D^x(\alpha,d)$ is realized
by a unique $H$--path, so each argument, when it attains the minimum, gives a
unique $\alpha$--$q$ path. Two distinct shortest $\alpha$--$q$ paths would force
$D^x(\alpha,c)+s=D^x(\alpha,d)+(x_f-s)$ --- which is exactly (G) for the edge
$f$, offset $s$, and the \emph{branch} target $q':=\alpha$ (the special case
(I)), and is therefore excluded. This proves the Claim. \hfill$\diamond$

\smallskip
Now take arbitrary vertices $p,q$ and suppose, for contradiction, that there are
two distinct shortest $p$--$q$ paths $P_1,P_2$. If both $p,q$ are branch this
contradicts (U). Otherwise we may assume $p$ is internal, on $e=ab$ at offset
$r\in\{1,\dots,x_e-1\}$; since $p$ has degree $2$, each $P_i$ leaves $p$ towards
$a$ or towards $b$.

\emph{Case A: $P_1,P_2$ leave $p$ in the same direction}, say towards $a$. The
vertices of $e$ strictly between $p$ and $a$ have degree $2$, so both paths
contain the segment $p\to a$; truncating it yields two distinct shortest
$a$--$q$ paths, contradicting the Claim (branch source $a$).

\emph{Case B: $P_1,P_2$ leave $p$ in opposite directions.} Then
$P_1=p\to a\to\cdots\to q$ has length $r+d_{G'}(a,q)$ and
$P_2=p\to b\to\cdots\to q$ has length $(x_e-r)+d_{G'}(b,q)$, and both equal
$d_{G'}(p,q)$, so
\begin{equation}\label{eq:caseB}
 r+d_{G'}(a,q)=(x_e-r)+d_{G'}(b,q).
\end{equation}
If $q$ does \emph{not} lie on $e$, then~\eqref{eq:caseB} is exactly a violation
of (G) (edge $e$, offset $r$, off-edge target $q$), contradiction.

It remains to treat $q$ \emph{on} $e$, say at offset $r'$ with (without loss of
generality) $r<r'<x_e$. Leaving $p$ towards $b$ one meets $q$ \emph{before} $b$
(all intermediate vertices have degree $2$), so the only $p$--$q$ path in that
direction is the on-edge segment, of length $r'-r$; hence $P_1$, the path leaving
$p$ towards $a$, must avoid that segment, i.e.\ $P_1=p\to a\to\pi\to b\to q$ where
$\pi$ is an $H$--path from $a$ to $b$ \emph{not using} $e$, of weight
$R:=\sum_{f\in\pi}x_f$. Equality of the two lengths in~\eqref{eq:caseB} forces
\begin{equation}\label{eq:Rrigid}
 r+R+(x_e-r')=r'-r,\qquad\text{i.e.}\qquad R=2(r'-r)-x_e .
\end{equation}
From $0<r<r'<x_e$ we read off $0<R<x_e$, so the edge $e$ is \emph{redundant}: its
single edge (weight $x_e$) is not a minimum-weight $a$--$b$ path. Since every
prefix of the geodesic $P_1$ is a geodesic, $a\to\pi\to b$ is a minimum-weight
$a$--$b$ path, so $D^x(a,b)=R$; by (U) the minimum path is unique, namely $\pi$.
Finally, from~\eqref{eq:Rrigid}, $x_e+R=2(r'-r)$ is \emph{even}.

Thus $e$ is a redundant edge with $x_e+R$ even. By Lemma~\ref{lem:redundant}
(whose hypothesis (U) we have assumed), (G) fails --- contradicting our
assumption that (G) holds. Hence Case~B with $q$ on $e$ cannot occur.

In every case we reached a contradiction, so the shortest $p$--$q$ path is
unique; $G'(x)$ is geodetic.
\end{proof}

\begin{remark}[Why the closure works]\label{rem:why}
Lemma~\ref{lem:redundant} is the heart of the matter. An on-edge double geodesic
is an antipodal pair on the even fundamental cycle $C=e\cup\pi$ of length
$x_e+R=2m$. The lemma observes that the \emph{same} cycle automatically carries an
\emph{off}-edge double geodesic: take any branch vertex $v$ interior to $\pi$; its
$C$--antipode lands at the interior offset $s=x_e+\beta-m$ of $e$, and the two
$C$--arcs from that point to $v$ both have length $m=d_{G'}(s\text{-point},v)$
because $s+\alpha=(x_e-s)+\beta=m$. The existence of an interior branch vertex $v$
on $\pi$ is guaranteed precisely by $\operatorname{girth}(H)=5$ (so $\pi$ has
$\ge 4$ edges). This is what makes the global condition (G) --- ranging over
branch targets reachable around long cycles --- already detect the on-edge
phenomenon, with no appeal to computation.
\end{remark}

\begin{remark}[The branch-only condition (I) is necessary but not
sufficient]\label{rem:counterex}
Take $x_e\equiv3$ for all edges except one edge $e_0$ with $x_{e_0}=5$. This
assignment satisfies (U) and the branch-only condition (I), yet $G'(x)$ is
\emph{not} geodetic: a direct shortest-path count finds pairs of vertices with
two shortest paths, all of them pairs of \emph{internal} vertices. The
obstruction is an opposite-direction tie
$r+d_{G'}(a,q)=(x_e-r)+d_{G'}(b,q)$ in which the target $q$ is itself internal,
invisible to (I), which only constrains \emph{branch} targets. Consequently any
characterization restricted to a fixed finite family of short cycles --- in
particular any purely pentagon/hexagon system --- is incomplete: condition (G) is
inherently global, referring to weighted shortest-path distances to all vertices
of $(H,x)$. (Theorem~\ref{thm:main} shows the corrected predicate
$(U)\wedge(G)$ is nevertheless exactly geodeticity.)
\end{remark}

\section{The solution set is infinite}

\begin{proposition}[Uniform odd scalings]\label{prop:infinite}
For every odd positive integer $c$, the uniform assignment $x_e\equiv c$ yields a
geodetic homeomorph $G'(x)$ of diameter $2c$; for every even $c\ge2$ it is not
geodetic. Hence there are infinitely many pairwise non-isomorphic geodetic
homeomorphs, with unbounded diameter.
\end{proposition}

\begin{proof}
With $x_e\equiv c$ the weighted distance equals $c$ times the unweighted
distance: $D^x(\alpha,\beta)=c\,n(\alpha,\beta)$ for branch vertices, where
$n(\cdot,\cdot)\in\{0,1,2\}$ is the unweighted $H$--distance. Since $H$ is
geodetic, weighted and unweighted minimum $H$--paths coincide and are unique,
giving (U). We verify (G) and then invoke Theorem~\ref{thm:main}.

First note that for uniform $c$ \emph{no edge is redundant}: an $a$--$b$ path
avoiding the edge $e=ab$ has, by girth $5$, at least $4$ edges, hence weight
$\ge 4c>c=x_e$; so $D^x(a,b)=x_e$ and $e$ itself is the unique minimum path.
In particular the on-edge case of (G) never arises.

Now consider an off-edge instance of (G): an edge $e=ab$, offset
$r\in\{1,\dots,c-1\}$, and a target $q$ off $e$. We must show
$r+d_{G'}(a,q)\neq(c-r)+d_{G'}(b,q)$, i.e.\
\begin{equation}\label{eq:tie}
 d_{G'}(a,q)-d_{G'}(b,q)\neq c-2r .
\end{equation}
If $q$ is a branch vertex: the left side is $c\,(n(a,q)-n(b,q))$, and since $a,b$
are adjacent $n(a,q)-n(b,q)\in\{-1,0,1\}$, so $d_{G'}(a,q)-d_{G'}(b,q)\in
\{-c,0,c\}$; but for $r\in\{1,\dots,c-1\}$ the value $c-2r$ is \emph{odd}
(hence nonzero, $c$ being odd) and $|c-2r|\le c-2<c$, so it cannot equal any of
$-c,0,c$, and~\eqref{eq:tie} holds. (This is condition (I).)

If $q$ is internal, on $f=\{c_0,d_0\}$ at offset $t\in\{1,\dots,c-1\}$ from
$c_0$, then by~\eqref{eq:internaldist}
\[
 d_{G'}(a,q)=\min\big(c\,n(a,c_0)+t,\ c\,n(a,d_0)+(c-t)\big),
\]
and similarly for $b$. Reducing modulo $c$, each of $d_{G'}(a,q),d_{G'}(b,q)$ is
congruent to $t$ or to $-t$. Hence
$d_{G'}(a,q)-d_{G'}(b,q)\equiv 0,\ 2t,\ \text{or}\ -2t \pmod c$. If~\eqref{eq:tie}
\emph{failed}, the left side would equal $c-2r\equiv -2r\pmod c$, forcing
$2r\equiv 0,\ \pm 2t\pmod c$; as $c$ is odd, $2$ is invertible mod $c$, so
$r\equiv 0,\ \pm t\pmod c$, i.e.\ (since $r,t\in\{1,\dots,c-1\}$) $r=t$ or
$r=c-t$.

Suppose $r=t$ (the case $r=c-t$ is symmetric, exchanging the roles of the two
ends of $f$). Write the realizing sides: $d_{G'}(a,q)=c\,n_a+t_a$ and
$d_{G'}(b,q)=c\,n_b+t_b$ where each $(n_\bullet,t_\bullet)$ is $(n(\cdot,c_0),t)$
or $(n(\cdot,d_0),c-t)$. Their difference must equal $c-2t$. If $t_a=t_b$ (both
sides enter $f$ at the same end), the difference is $c(n_a-n_b)$, and
$c(n_a-n_b)=c-2t$ forces $2t=c(1-(n_a-n_b))$, a positive multiple of $c$ in
$\{2,\dots,2c-2\}$, hence $2t=c$, impossible for $c$ odd. If $t_a\neq t_b$ (the
two sides enter $f$ at opposite ends), the difference is
$c(n_a-n_b)\pm(t-(c-t))=c(n_a-n_b)\pm(2t-c)$, and setting this to $c-2t$ gives
$c\,(n_a-n_b)\mp(c-2t)=c-2t$, i.e.\ $c(n_a-n_b)=(1\pm1)(c-2t)\in\{0,\,2(c-2t)\}$.
The value $0$ forces $n_a=n_b$ and $c-2t=0$ (impossible, $c$ odd); the value
$2(c-2t)$ forces $c\mid 2(c-2t)$, hence $c\mid 4t$, hence (as $\gcd(c,4)=1$)
$c\mid t$, impossible for $0<t<c$. In every sub-case~\eqref{eq:tie} holds.

Thus (G) holds, and by Theorem~\ref{thm:main} $G'(x)$ is geodetic. Its diameter
is $c\cdot\operatorname{diam}(H)=2c$, so distinct odd $c$ give non-isomorphic
graphs.

For even $c\ge2$, take any edge $e=ab$ and a branch vertex $t$ with
$n(t,b)=n(t,a)+1$ (an out-neighbour direction); then
$r+d_{G'}(a,t)=(c-r)+d_{G'}(b,t)$ at $r=c/2\in\{1,\dots,c-1\}$, an explicit
opposite-direction tie with branch target, so (I) and hence geodeticity fail.
\end{proof}

\begin{corollary}\label{cor:illposed}
The task ``produce a complete list and exact \emph{count} of all geodetic graphs
of diameter $d\ge2$ homeomorphic to $H$'' has no finite answer: by
Proposition~\ref{prop:infinite} there are infinitely many, even up to
isomorphism. A finite enumeration is meaningful only after a finiteness
normalization, e.g.\ bounding $\max_e x_e\le D$ (equivalently bounding the
subdivision lengths, hence the diameter).
\end{corollary}

\section{Algorithm}

Fix a bound $D$ and enumerate exactly the assignments with $\max_e x_e\le D$
for which $G'(x)$ is geodetic. The simplest provably correct test applies the
definition of geodeticity to $G'(x)$ directly; the characterization
Theorem~\ref{thm:main} provides a faster equivalent test on $(H,x)$.

\begin{algorithm}[H]
\caption{Enumerate geodetic homeomorphs of $H$ with $\max_e x_e\le D$}
\begin{algorithmic}[1]
\Require $H$ ($50$ vertices, $175$ edges); a bound $D\ge1$.
\Ensure all $x\in\{1,\dots,D\}^{E(H)}$ with $G'(x)$ geodetic.
\State $\mathcal S\gets\emptyset$.
\For{each $x\in\{1,\dots,D\}^{E(H)}$ (generated by backtracking with the
necessary-condition pruning below)}
  \State Build $G'(x)$ (it has $50+\sum_e(x_e-1)\le 50+175(D-1)$ vertices).
  \State \textbf{Geodeticity test:} from every vertex run a BFS that counts
  shortest paths; \emph{reject} as soon as some vertex is reached by more than
  one shortest path. \Comment{equivalently: test (U) by a shortest-path-count
  Dijkstra on $(H,x)$, then test (G) by its arithmetic restatement, using
  \eqref{eq:internaldist} for the internal distances.}
  \State If no rejection occurs, add $x$ to $\mathcal S$.
\EndFor
\State \Return $\mathcal S$.
\end{algorithmic}
\end{algorithm}

\begin{proposition}[Correctness]
The algorithm returns exactly $\{x\in\{1,\dots,D\}^{E(H)}:G'(x)\text{ geodetic}\}$.
\end{proposition}
\begin{proof}
The BFS shortest-path-count test accepts $G'(x)$ iff every ordered pair of
vertices has a unique shortest path, i.e.\ iff $G'(x)$ is geodetic --- this is
the definition. The alternative $(U)\wedge(G)$ test is equivalent by
Theorem~\ref{thm:main}. The pruning (below) discards only non-geodetic $x$, since
the conditions used are necessary. Hence the accepted set is exactly the geodetic
assignments with $\max_e x_e\le D$.
\end{proof}

\begin{proposition}[Termination and per-candidate cost]
The algorithm halts; the search space has $\le D^{175}$ leaves. For each
candidate the geodeticity test costs $O(|V'|\cdot|E'|)$ with
$|V'|,|E'|=O(175\,D)$; equivalently the $(U)\wedge(G)$ test costs
$O\!\big(50\cdot|E|\log|V|\big)$ for the $50$ Dijkstra runs plus
$O(|V|\cdot|E|\cdot D)$ for the (G) scan, polynomial per candidate.
\end{proposition}

\noindent\textbf{Pruning.} Any geodetic $x$ satisfies the following \emph{necessary}
local conditions, usable to prune the backtracking cheaply:
(P) every pentagon has odd total weight; (H) on every hexagon the two arcs of
each $1{+}5$ and $2{+}4$ split have distinct weights; and more generally the
branch-only condition (I). These are special cases of $(U)\wedge(G)$ for branch
pairs and branch targets and may be checked incrementally, but they may
\emph{never} replace the full geodeticity test (Remark~\ref{rem:counterex}).

\begin{remark}[Honest assessment of practicality]
The per-candidate test is polynomial, but the candidate space $D^{175}$ is
astronomically large, so exhaustive enumeration is infeasible beyond very small
$D$. The necessary local conditions give strong but \emph{insufficient} pruning.
We do not claim a polynomial-time enumeration in $D$; we provide a correct
decision procedure, a correct global characterization, and a sound pruning
framework. For $D=1$ only $x\equiv1$ survives, returning $H$ itself; among
constant assignments exactly the odd $c\le D$ survive
(Proposition~\ref{prop:infinite}).
\end{remark}

\section{On the figure ``$5376$'' and on automorphisms}

The problem statement asserts that the geodetic conditions form a system of
$5376$ equations over the $5$- and $6$-cycles. Remark~\ref{rem:counterex} shows
that no such purely local (pentagon/hexagon) system can be complete: the decisive
condition (G) is global, involving weighted shortest-path distances to internal
vertices on \emph{arbitrarily long} cycles. We were unable to reproduce $5376$
from natural local bookkeeping ($1260$ pentagons; $1260+5250$ pentagons and
hexagons; etc., all differ) and record this discrepancy honestly. The
automorphism group $\operatorname{Aut}(H)\cong\mathrm P\Sigma\mathrm U(3,5)$ of
order $252000$ acts on the solution set; same-orbit assignments yield isomorphic
homeomorphs, so for any finite normalization the number of isomorphism classes
follows from the fixed-point counts by Burnside's lemma.

\section{Summary of what is established}

\begin{enumerate}
\item \textbf{(Characterization, Theorem~\ref{thm:main}.)} $G'(x)$ is geodetic
iff (U) every branch pair has a unique minimum-weight $H$--path, and (G) there is
no opposite-direction tie $r+d_{G'}(a,q)=(x_e-r)+d_{G'}(b,q)$ for an internal
point of an edge $e=ab$ and \emph{any} vertex $q$ off $e$. \emph{Both necessity
and sufficiency are proved in full.} The previously open degenerate sub-case (two
internal points of the same edge) is closed in closed form by
Lemma~\ref{lem:redundant}: a redundant edge with even $x_e+R$ forces, via an
interior branch vertex of the closing path $\pi$ (which exists because
$\operatorname{girth}(H)=5$), an off-edge violation of (G).
\item \textbf{(Local conditions are insufficient, Remark~\ref{rem:counterex}.)}
The branch-only condition (I) --- and a fortiori pentagon parity and hexagon
non-tie --- is necessary but not sufficient; an explicit counterexample (uniform
$3$ with one edge of length $5$) satisfies $(U)\wedge(I)$ yet is non-geodetic,
the failure occurring at internal--internal pairs.
\item \textbf{(Infiniteness, Proposition~\ref{prop:infinite},
Corollary~\ref{cor:illposed}.)} Every uniform odd assignment is geodetic of
diameter $2c$ (every even one is not), proved in full; hence there are infinitely
many geodetic homeomorphs and no finite exact count without a normalization.
\item \textbf{(Algorithm.)} A correct decision procedure builds $G'(x)$ and
tests geodeticity by BFS shortest-path counts, equivalently by $(U)\wedge(G)$ on
$(H,x)$; it is polynomial per candidate. Necessary local pruning is sound but
cannot replace the global test. No polynomial-time enumeration in $D$ is claimed.
\end{enumerate}

\noindent\textbf{Reproducibility.} The numerical sanity checks behind this
write-up are elementary BFS/Dijkstra computations on the actual Hoffman--Singleton
graph (Robertson pentagon/pentagram model, confirmed $50$ vertices, $7$--regular,
$175$ edges, triangle-free, diameter $2$). They corroborate but are not relied
upon by the proofs above, which are self-contained. Specifically, building $G'(x)$
and testing geodeticity by per-vertex shortest-path-count BFS confirms:
(i) the equivalence $(U)\wedge(G)\Leftrightarrow$ geodetic over $400$ diverse
random subdivisions of $H$ (regimes $x_e\in\{1,\dots,4\}$, the tie-prone
$x_e\in\{1,2\}$, and uniform-base-with-perturbations), with \emph{zero}
mismatches; (ii) the counterexample of Remark~\ref{rem:counterex} (uniform $3$
with one edge of length $5$) satisfies $(U)$ and the branch-only $(I)$ yet is
non-geodetic, the obstruction realized by $(G)$ at an \emph{internal} target;
and (iii) the uniform family of Proposition~\ref{prop:infinite}: $x_e\equiv c$ is
geodetic for $c\in\{1,3,5,7\}$ and non-geodetic for $c\in\{2,4,6\}$.

\end{document}
